% from intro-fr.tex
\oldpage[I]{8}
de vue de la dimension (chap.~IV), qui se transcrit sans changement notable dans le cadre des schémas.
Il convient de noter par ailleurs que la notion de <<~schémas d'anneaux locaux~>>, introduite par Chevalley, se prête naturellement à une extension de la Géométrie algébrique (n'ayant pas cependant toute la souplesse et la généralité que nous entendons lui donner ici); pour les rapports entre cette notion et notre théorie, voir chapitre premier, \S~8.
Une telle extension a été développée par M.~Nagata dans une série de mémoires~\cite{I-9} contenant de nombreux résultats spéciaux concernant la Géométrie algébrique sur les anneaux de Dedekind\footnote{Parmi les travaux qui se rapprochent de notre point de vue en Géométrie algébrique, signalons l'important travail de E.~Kähler~\cite{I-22}, et une Note récente de Chow et Igusa~\cite{I-3}, qui reprennent dans le cadre de la théorie de Nagata-Chevalley certains résultats de (FAC) et donnent aussi une formule de Künneth.}.

\bigskip
\begin{center}*\quad*\quad*\end{center}
\bigskip

Enfin, il va sans dire qu'un livre sur la Géométrie algébrique, et surtout un livre portant sur les fondements, est nécessairement influencé, ne serait-ce que par personnes interposées, par des mathématiciens tels que O.~Zariski et A.~Weil.
En particulier, la \emph{Théorie des fonctions holomorphes} de Zariski~\cite{I-20}, convenablement assouplie grâce aux méthodes cohomologiques et complétée par un théorème d'existence (chap.~III, \S\S~4 et 5) est (avec la technique de descente exposée au chap.~VI) un des principaux outils employés dans ce Traité, et nous semble un des plus puissants dont on dispose en Géométrie algébrique.

La technique générale dans laquelle elle s'insère peut être esquissée de la façon suivante (un exemple typique en sera fourni au chap.~IX, dans l'étude du groupe fondamental).
On a un morphisme propre (chap.~II) $f:X\to Y$ d'une variété algébrique dans une autre (plus généralement, d'un schéma dans un autre) qu'on veut étudier au voisinage d'un point $y\in Y$, en vue de résoudre un problème P relatif à un voisinage de $y$.
On opère par étapes successives~:
\begin{enumerate}
  \item[$1^\circ$] On peut supposer $Y$ affine, de sorte que $X$ devient un schéma défini sur l'anneau affine $A$ de $Y$, et on peut même remplacer $A$ par l'anneau local de $y$.
  Cette réduction est toujours facile en pratique (chap.~V) et nous ramène au cas où $A$ est un anneau \emph{local}.
  \item[$2^\circ$] On étudie le problème envisagé lorsque $A$ est un anneau local \emph{artinien}.
  Pour qu'il garde effectivement un sens lorsque $A$ n'est pas supposé intègre, il y a lieu parfois de reformuler le problème P, et il apparaît que l'on obtient souvent ainsi une meilleure compréhension du problème, de nature <<~infinitésimale~>> à ce stade.
  \item[$3^\circ$] La théorie des schémas formels (chap.~III, \S\S~3, 4 et 5) permet de passer du cas d'un anneau artinien au cas d'un \emph{anneau local complet}.
  \item[$4^\circ$] Enfin, si $A$ est un anneau local quelconque, la considération de <<~sections multiformes~>> sur des schémas convenables sur $X$, approchant une section <<~formelle~>> donnée (chap.~IV) permettra souvent de passer d'un résultat connu pour le schéma

  